%--------------------------------------------------
%  Diagrams of decorated unit-interval graphs
%--------------------------------------------------
%
% A graph is drawn as in [AS22]: vertex v is the diagonal box
% [v-1,v]x[v-1,v], the natural edge uv (u<v) is the cell [u-1,u]x[v-1,v],
% and the Schröder path in the steps n, e, d bounds the natural edges from
% the left.  A diagonal step crosses a strict corner.
%
%   \LLTpic[<tikz options>]{<path>}{<overlays>}
%       A complete diagram, e.g. \LLTpic{nnnddeneee}{\LLTasc{1/3,2/3}}.
%   \LLTdiagram{<path>}      inside a tikzpicture with style llt
%   \LLTasc{u/v,...}         ascending edges (arrows)
%   \LLTvertex{v,...}{color} recolor diagonal boxes
%   \LLTshade{u/v,...}{color} shade cells
%   \LLTmark{u/v,...}{text}  put a symbol in cells
%   \LLTsteps[style]{x}{y}{path}  draw a path from (x,y), for general
%       pictures; combine with the styles llt arbitrary, llt diagonal and
%       llt bounce.
%   \LLTframe{w}{h}          bounding box and diagonal of a general picture
%
%   \EdgePic[spacing]{k}{labels}{edges}  k vertices on a line, named by the comma
%       separated labels, on a dotted line standing for the other vertices and
%       edges; edges are i/j/s (strict), i/j/a (ascending) or i/j/c (a
%       retained cover, gray), drawn as arcs above, or i/j/b (ascending,
%       drawn below).
%
% Options: llt/cell (box size, default 1.1em), llt/label font,
% llt/labels (true/false), llt/first (label of the first vertex, default 1).

\pgfdeclarelayer{lltback}
\pgfsetlayers{lltback,main}

\makeatletter
\newif\ifllt@labels
\pgfkeys{/llt/.cd,
  cell/.store in=\llt@cell, cell=1.1em,
  label font/.store in=\llt@labelfont, label font=\scriptsize,
  labels/.is if=llt@labels, labels=true,
  first/.store in=\llt@first, first=1,
  vertex fill/.store in=\llt@vfill, vertex fill=yellow}
\newcommand{\lltset}[1]{\pgfkeys{/llt/.cd,#1}}

\tikzset{
  llt/.style={/llt/.cd,#1,/tikz/.cd,x=\llt@cell,y=\llt@cell},
  llt/.default={},
  llt grid/.style={draw=black!55,line width=0.3pt},
  llt path/.style={line width=1.1pt,line cap=round,line join=round},
  llt arbitrary/.style={line width=1.1pt,dotted,line cap=round},
  llt diagonal/.style={black!45,line width=0.4pt},
  llt bounce/.style={red!75!black,line width=0.9pt,densely dashed,
    -{Stealth[length=1.6mm]}},
  llt asc/.style={-{Stealth[length=1mm]},line width=0.5pt},
  llt cover/.style={-{Stealth[length=2.2mm]},line width=1.4pt,black!45},
}

\newcount\llt@x
\newcount\llt@y
\newcount\llt@i

% \LLTsteps[style]{x}{y}{word}
\newcommand{\LLTsteps}[4][llt path]{%
  \begingroup
  \pgfmathsetmacro\llt@px{#2}\pgfmathsetmacro\llt@py{#3}%
  \edef\llt@path{(\llt@px,\llt@py)}%
  \@tfor\llt@s:=#4\do{%
    \if n\llt@s\pgfmathsetmacro\llt@py{\llt@py+1}\fi
    \if e\llt@s\pgfmathsetmacro\llt@px{\llt@px+1}\fi
    \if d\llt@s\pgfmathsetmacro\llt@px{\llt@px+1}%
      \pgfmathsetmacro\llt@py{\llt@py+1}\fi
    \edef\llt@path{\llt@path--(\llt@px,\llt@py)}}%
  \draw[#1] \llt@path;%
  \endgroup}

% One row: the path is at x, the row is y+1.
\newcommand{\llt@row}{%
  \llt@i=\llt@x
  \loop\ifnum\llt@i<\llt@y
    \draw[llt grid] (\the\llt@i,\the\llt@y) rectangle +(1,1);%
    \advance\llt@i by 1
  \repeat
  \begin{pgfonlayer}{lltback}
    \fill[\llt@vfill] (\the\llt@y,\the\llt@y) rectangle +(1,1);%
  \end{pgfonlayer}%
  \draw[llt grid] (\the\llt@y,\the\llt@y) rectangle +(1,1);%
  \ifllt@labels
    \node[font=\llt@labelfont,inner sep=0pt]
      at (\the\llt@y+0.5,\the\llt@y+0.5) {$\the\numexpr\llt@y+\llt@first\relax$};%
  \fi}

\newcommand{\LLTdiagram}[1]{%
  \begin{scope}[shift={(\llt@first-1,\llt@first-1)}]
  \begingroup
  \llt@x=0 \llt@y=0
  \@tfor\llt@s:=#1\do{%
    \if n\llt@s\llt@row\advance\llt@y by 1 \fi
    \if d\llt@s\llt@row\advance\llt@x by 1 \advance\llt@y by 1 \fi
    \if e\llt@s\advance\llt@x by 1 \fi}%
  \endgroup
  \LLTsteps{0}{0}{#1}
  \end{scope}}

\newcommand{\LLTframe}[2]{%
  \path (-0.4,-0.4) rectangle (#1+0.4,#2+0.4);%
  \draw[llt diagonal] (-0.4,-0.4)--(#1+0.4,#1+0.4);}

\newcommand{\EdgePic}[4][2em]{%
  \begin{tikzpicture}[x=#1,y=#1,baseline={([yshift=-0.6ex]v1.center)},
      every node/.append style={minimum size=1.05em,font=\scriptsize}]
  \draw[llt arbitrary,line width=0.6pt] (0.4,0)--(#2+0.6,0);
  \foreach \l [count=\i] in {#3}{\node[vtx,fill=white] (v\i) at (\i,0) {$\l$};}
  \foreach \i/\j/\s in {#4}{%
    \pgfmathsetmacro\llt@ang{35+12*(\j-\i)}%
    \def\llt@st{sel}%
    \if\s s\def\llt@st{strictedge}\fi
    \if\s c\def\llt@st{llt cover}\fi
    \if\s b\draw[sel] (v\i) to[out=-\llt@ang,in=180+\llt@ang] (v\j);%
    \else\draw[\llt@st] (v\i) to[out=\llt@ang,in=180-\llt@ang] (v\j);\fi}
  \end{tikzpicture}}

\newcommand{\LLTasc}[1]{%
  \foreach \u/\v in {#1}{\draw[llt asc] (\u-0.82,\v-0.5)--(\u-0.18,\v-0.5);}}

\newcommand{\LLTvertex}[2]{%
  \foreach \v in {#1}{%
    \begin{pgfonlayer}{lltback}
      \fill[#2] (\v-1,\v-1) rectangle (\v,\v);
    \end{pgfonlayer}}}

\newcommand{\LLTshade}[2]{%
  \foreach \u/\v in {#1}{%
    \begin{pgfonlayer}{lltback}
      \fill[#2] (\u-1,\v-1) rectangle (\u,\v);
    \end{pgfonlayer}}}

\newcommand{\LLTmark}[2]{%
  \foreach \u/\v in {#1}{%
    \node[font=\llt@labelfont,inner sep=0pt] at (\u-0.5,\v-0.5) {#2};}}

\newcommand{\LLTpic}[3][]{%
  \begin{tikzpicture}[llt={#1}]\LLTdiagram{#2}#3\end{tikzpicture}}
\makeatother
